\begin{figure}[ht]
\centering
\begin{tikzpicture}[x=1cm, y=1cm, font=\small]

% Cauchy condensation for a monotone-decreasing term a_n. The bars are the
% terms 1/n across dyadic blocks n in [2^k, 2^{k+1}); the horizontal caps
% show each block replaced by its largest term (upper bound) and its
% smallest term (lower bound), so the block sum is trapped between
% 2^k * a_{2^{k+1}} and 2^k * a_{2^k}. The two condensed bounds sandwich
% the original sum, which is why sum a_n and sum 2^k a_{2^k} share a verdict.
% x-axis is n on a log2 spacing; y is term magnitude, drawn to scale.

\draw[->, thick] (-0.2, 0) -- (10.4, 0) node[right] {index $n$ (dyadic blocks)};
\draw[->, thick] (0, -0.2) -- (0, 4.2) node[above] {term $a_{n}=1/n$};

% draw term bars for a few dyadic blocks; heights scaled by 4*(1/n) capped
% block k=0: n=1
\draw[fill=engblue!25, draw=engblue] (0.30,0) rectangle (0.70, 4.000);
% block k=1: n=2,3  (heights 4/2=2.0, 4/3=1.333)
\draw[fill=engblue!25, draw=engblue] (1.10,0) rectangle (1.50, 2.000);
\draw[fill=engblue!25, draw=engblue] (1.55,0) rectangle (1.95, 1.333);
% block k=2: n=4..7 (heights 1.0,0.8,0.667,0.571)
\draw[fill=engblue!25, draw=engblue] (2.60,0) rectangle (2.95, 1.000);
\draw[fill=engblue!25, draw=engblue] (2.98,0) rectangle (3.33, 0.800);
\draw[fill=engblue!25, draw=engblue] (3.36,0) rectangle (3.71, 0.667);
\draw[fill=engblue!25, draw=engblue] (3.74,0) rectangle (4.09, 0.571);
% block k=3: n=8..15 (first height 0.5, last 4/15=0.267); draw as a band
\foreach \i/\h in {0/0.500,1/0.444,2/0.400,3/0.364,4/0.333,5/0.308,6/0.286,7/0.267} {
  \pgfmathsetmacro\xa{5.10 + \i*0.30}
  \pgfmathsetmacro\xb{\xa + 0.26}
  \draw[fill=engblue!25, draw=engblue] (\xa,0) rectangle (\xb, \h);
}

% upper-bound caps: each block replaced by its largest term (leftmost height)
\draw[very thick, engred] (1.10, 2.000) -- (1.95, 2.000);
\draw[very thick, engred] (2.60, 1.000) -- (4.09, 1.000);
\draw[very thick, engred] (5.10, 0.500) -- (7.46, 0.500);
\node[engred, anchor=south, font=\footnotesize] at (6.28, 0.52) {upper cap $2^{k}a_{2^{k}}$};

% lower-bound caps: each block replaced by its smallest term (rightmost height)
\draw[very thick, engorange, dashed] (2.60, 0.571) -- (4.09, 0.571);
\draw[very thick, engorange, dashed] (5.10, 0.267) -- (7.46, 0.267);
\node[engorange, anchor=north, font=\footnotesize] at (6.28, 0.25) {lower cap $2^{k}a_{2^{k+1}}$};

% block labels on x-axis
\foreach \x/\lab in {0.50/{$1$}, 1.52/{$2..3$}, 3.34/{$4..7$}, 6.28/{$8..15$}} {
  \node[anchor=north, font=\footnotesize] at (\x, -0.05) {\lab};
}

\end{tikzpicture}
\caption[Cauchy condensation of a monotone series]{Cauchy condensation
of a monotone-decreasing series, shown on $a_{n}=1/n$. The terms are
grouped into dyadic blocks $n\in[2^{k},2^{k+1})$, each holding $2^{k}$
terms. Replacing every term in a block by the block's largest term (solid
cap) gives the upper bound $2^{k}a_{2^{k}}$ on the block sum; replacing by
the smallest term (dashed cap) gives the lower bound $2^{k}a_{2^{k+1}}$.
The condensed series $\sum 2^{k}a_{2^{k}}$ sits between one copy and one
half-copy of the original tail, so the two series converge or diverge
together. For $a_{n}=1/n$ the condensed series is $\sum 2^{k}\cdot
2^{-k}=\sum 1$, plainly divergent, which is the harmonic divergence read
off one dyadic bound rather than the Oresme grouping.}
\label{fig:vol02:ch04:cauchy-condensation}
\end{figure}
